% Resume: Arumalla Mohit Krishna
% Compile with pdfLaTeX.

\documentclass[letterpaper,10pt]{article}

\usepackage[margin=0.5in]{geometry}
\usepackage[T1]{fontenc}
\usepackage{lmodern}
\usepackage{microtype}
\usepackage{titlesec}
\usepackage{enumitem}
\usepackage[hidelinks]{hyperref}
\usepackage{tabularx}
\input{glyphtounicode}

\pdfgentounicode=1
\pagestyle{empty}
\urlstyle{same}
\raggedright
\setlength{\parindent}{0pt}
\setlength{\tabcolsep}{0pt}

\titleformat{\section}
  {\large\scshape}{}{0em}{}[\titlerule]

\titlespacing*{\section}{0pt}{6pt}{3pt}

\newcommand{\entry}[2]{%
  \noindent
  \begin{tabularx}{\linewidth}
    {@{}>{\raggedright\arraybackslash}X@{\hspace{6pt}}r@{}}
    \textbf{#1} & #2
  \end{tabularx}\par
}

\newcommand{\bullets}{
  \begin{itemize}[
    leftmargin=1.2em,
    topsep=2pt,
    itemsep=1pt,
    parsep=0pt,
    partopsep=0pt
  ]
}

\newcommand{\finishbullets}{
  \end{itemize}
}

\begin{document}

%----------HEADING----------
\begin{center}
    {\LARGE\bfseries Arumalla Mohit Krishna}\\[3pt]
    +1 (631) 352-8643 $|$
    \href{mailto:mohitkrishna.arumalla@stonybrook.edu}
    {mohitkrishna.arumalla@stonybrook.edu} $|$
    \href{https://linkedin.com/in/mohit-krishna-arumalla}
    {LinkedIn}
\end{center}
\vspace{-8pt}

%-----------EDUCATION-----------
\section{Education}

\entry
  {\textbf{Stony Brook University}}
  {\textbf{Aug 2025 -- Dec 2026}}

M.S. in Computer Science \hfill New York, USA\\
\textit{
Coursework: Natural Language Processing,
Introduction to Computer Vision, Vision Language Modeling
}
\vspace{3pt}

\entry
  {\textbf{Birla Institute of Technology and Science, Pilani}}
  {\textbf{Aug 2019 -- 2023}}

B.E. in Electronics \& Communication; GPA: 8.02/10 \hfill India

%-----------PROFESSIONAL SUMMARY-----------
\section{Professional Summary}

\textbf{Software engineer} with 2+ years of experience building \textbf{distributed services},
\textbf{AI-powered semantic search}, and real-time data pipelines, pursuing an M.S. in Computer
Science at Stony Brook University. Hands-on experience designing, deploying, and integrating
\textbf{AI systems} including large language model based retrieval and inference services. Skilled in
Python, Go, PyTorch, containerization, and cloud-native infrastructure with a strong foundation in
delivering \textbf{production-grade AI software solutions}.

%-----------TECHNICAL SKILLS-----------
\section{Technical Skills}

\textbf{Languages:}
Python, Go, C++, SQL, Java\\

\textbf{AI \& ML Technologies:}
PyTorch, TensorFlow/Keras, Transformers, mBART, BGE, Vision Transformers, Self-Supervised Learning, LoRA, Visual Prompt Tuning, Text Embeddings, Vector Search, Fine-tuning, Structured Outputs, LLM Function Calling, Multi-Agent Orchestration\\

\textbf{Data \& Systems:}
gRPC, Kafka, Apache Flink, Flink SQL, Debezium/CDC,
PostgreSQL, Redis, Elasticsearch, Airflow, BigQuery\\

\textbf{Infrastructure:}
NVIDIA Triton, AWS, Docker, Kubernetes, Linux, Git,
CI/CD, Istio/Envoy, New Relic, Monitoring, Tracing, AI Observability

%-----------EXPERIENCE-----------
\section{Experience}

\entry
  {\textbf{Software Engineer II} $|$ \textbf{Infra.Market}, Bangalore}
  {\textbf{Apr 2024 -- Aug 2025}}

\bullets

  \item Built \textbf{semantic product search} using Amazon Titan text embeddings
  and Elasticsearch vector search, enabling \textbf{natural-language product discovery}
  and conversational query handling.

  \item Engineered a \textbf{CDC pipeline} using Flink, Kafka, Debezium,
  PostgreSQL, RocksDB, and Elasticsearch; reduced data synchronization latency
  from 10--15 minutes to sub-second with exactly-once semantics, improving data reliability.

  \item Deployed \textbf{Istio/Envoy service mesh} across 15+ microservices and integrated
  New Relic distributed tracing to enhance production monitoring, debugging, and reliability.

\finishbullets

\entry
  {\textbf{Software Engineer I} $|$ \textbf{Infra.Market}, Chennai}
  {\textbf{Jul 2023 -- Mar 2024}}

\bullets

  \item Designed and maintained \textbf{Search and Location microservices} with gRPC, serving over 500K daily requests for product discovery and geolocation with high reliability.

  \item Diagnosed and fixed PostgreSQL connection leaks using Gorm ORM, reducing
  connection-pool exhaustion by 85\%; validated fixes with load testing at 10K requests/second.

  \item Developed Airflow pipelines for BigQuery-to-Elasticsearch data migrations
  optimized for query performance, dynamic filtering, and automated monitoring to maintain index freshness.

\finishbullets

%-----------PROJECTS-----------
\section{Projects}

\entry
  {GPU-Optimized Personalized Product Search}
  {Aug 2026 -- Oct 2026}

\textit{Python, PyTorch, BGE, NVIDIA Triton, Go, gRPC}

\bullets

  \item Developed two-tower retrieval system using
  \texttt{BAAI/bge-small-en-v1.5} query/product embeddings
  (384 dimensions) with a trainable user tower combining query
  embeddings and user interaction history.

  \item Built an inference pipeline featuring precomputed product embeddings,
  approximate nearest-neighbor retrieval, and a Go/gRPC gateway;
  integrated NVIDIA Triton for GPU-accelerated query encoder serving.

  \item Designed benchmarking frameworks for batching and concurrency,
  measuring latency (p95/p99), throughput, GPU utilization, and cost efficiency;
  evaluated recall and ranking metrics while testing system overload recovery.

\finishbullets

\entry
  {Attention-Based Hallucination Detection in NMT}
  {Apr 2026 -- May 2026}

\textit{Python, PyTorch, mBART, Transformers, scikit-learn}

\bullets

  \item Built English-German hallucination-detection pipeline leveraging
  teacher-forced mBART attention features using WMT22/WMT23 quality estimation data.

  \item Extracted cross-/self-attention metrics across 12 decoder layers and 16 heads,
  including entropy, top-k mass, self-to-cross grounding ratios, enabling detection of hallucinations.

  \item Trained multi-level classifiers achieving 0.712 F1 and 0.695 ROC-AUC on 50K examples;
  proposed hallucination-aware reranking integrating risk scores and attention signals.

\finishbullets

\entry
  {Prompt/LoRA for Self-Supervised Learning}
  {Ongoing}

\textit{Python, PyTorch, Vision Transformers}

\bullets

  \item Researching domain-specific visual prompting and \textbf{LoRA} methods within
  a self-supervised ViT framework to learn transferable image representations across domains.

  \item Evaluating backbone model performance with/without prompts and LoRA at inference,
  analyzing classification accuracy and parameter efficiency under domain shifts.

\finishbullets

\end{document}